\chapter{Referencial Teórico}\label{cap-referencial-teorico}


\section{O E-Museu}

O E-Museu é uma iniciativa digital dedicada à preservação da história da computação, voltada à consulta pública e à gestão do acervo. Seu propósito é registrar e disponibilizar informações sobre a evolução da informática, de equipamentos clássicos a marcos do desenvolvimento de hardware, de modo que o conhecimento histórico permaneça acessível mesmo quando os artefatos físicos se deterioram.

O sistema oferece catálogo público com fichas, imagens e textos, além de painel administrativo para curadores registrarem, validarem e organizarem as peças. Entre os conteúdos, reúne dados técnicos sobre processadores, placas-mãe, memórias, unidades de armazenamento e periféricos, com descrição e histórico de uso.

O projeto adota modelo colaborativo: a comunidade pode contribuir com peças, documentos e relatos, sujeitos a verificação antes de integrarem a coleção. Voltado a fins educacionais e culturais, opera como iniciativa sem fins lucrativos, com foco na divulgação da memória tecnológica.


\section{Refatoração de software}

A refatoração é um processo fundamental no ciclo de desenvolvimento de sistemas, pois busca melhorar a estrutura interna do código sem alterar seu comportamento externo. No E-Museu, essa prática se torna importante para garantir qualidade, legibilidade, extensibilidade e facilidade de manutenção, especialmente diante da evolução constante de requisitos técnicos e funcionais.

De acordo com \citeonline{refatoracao:2004}, refatorar significa modificar a organização do software de forma sistemática, visando eliminar complexidades desnecessárias, reduzir duplicações, aprimorar o design e facilitar a compreensão do código. Refatoração não é uma etapa isolada, mas um processo contínuo que acompanha o desenvolvimento, permitindo melhorias incrementais conforme o sistema cresce.

Além disso, práticas de refatoração auxiliam na prevenção da chamada ``dívida técnica'', evitando que o código se torne difícil de compreender, testar ou modificar. Essa preocupação é especialmente relevante em sistemas acadêmicos e colaborativos, nos quais novas funcionalidades podem surgir a partir de pesquisas, contribuições externas ou mudanças institucionais.

No contexto deste trabalho, a refatoração contempla a reorganização e a melhoria estrutural dos controladores da versão legada do E-Museu, visando padronizar responsabilidades, reduzir acoplamentos indevidos e tornar a lógica de controle mais clara, testável e alinhada às boas práticas de desenvolvimento de software.

\citeonline{refatoracao:2004} também enfatiza o papel dos testes automatizados no processo de refatoração, de modo que o comportamento do software permaneça correto após cada melhoria estrutural. Dessa forma, refatorar não apenas aprimora o design do código, mas fortalece a confiabilidade do sistema.

\section{Boas práticas de desenvolvimento}

Esta seção apresenta as principais boas práticas consideradas no desenvolvimento do E-Museu e como cada abordagem contribui para organização, qualidade e manutenção do sistema.


\subsection{\acrfull{mvc}}

O padrão \acrfull{mvc} estabelece a separação do software em três camadas principais: \textit{Model}, \textit{View} e \textit{Controller} \cite{mvc}. Essa divisão organiza a aplicação de forma que a lógica de negócios, a interface e o controle de requisições não se misturem, reduzindo acoplamento e facilitando a manutenção. No contexto do E-Museu, o \acrshort{mvc} estrutura o fluxo de exibição do acervo, a manipulação dos dados e o controle das requisições, tornando o sistema mais claro e modular para quem der continuidade ao código.


\subsection{\textit{Clean Code}}

O princípio de \textit{Clean Code} orienta a escrita de códigos simples, legíveis, coesos e de fácil entendimento, priorizando qualidade sobre complexidade desnecessária \cite{cleancode:2009}. Sua aplicação incentiva boas práticas como nomes descritivos, funções curtas, ausência de códigos duplicados, clareza semântica e utilização consistente de padrões arquiteturais. Para o E-Museu, adotar \textit{Clean Code} contribui para a evolução do sistema, reduz custos de manutenção, facilita correções e permite que novos desenvolvedores entendam o projeto com menos esforço.

\subsection{\acrfull{ddd}}\label{sec:ref-ddd}

O \acrfull{ddd} propõe que o desenvolvimento de software seja guiado pelo domínio do negócio, priorizando o entendimento profundo do problema antes da solução tecnológica \cite{ddd2004}. Entre seus conceitos centrais destacam-se a \textit{linguagem ubíqua} (vocabulário compartilhado entre desenvolvedores e especialistas do museu), os \textit{contextos delimitados} (fronteiras em que um modelo vale de forma coerente) e a organização do código em camadas que refletem regras de negócio, e não apenas detalhes de infraestrutura.

No E-Museu, adotou-se uma aplicação pragmática dessa abordagem: o repositório foi reorganizado em pastas e \textit{namespaces} PHP alinhados a contextos de negócio (por exemplo, catálogo do acervo, taxonomia de etiquetas, colaboradores externos e identidade administrativa), com serviços, ações e requisições HTTP agrupados por domínio. Não se buscou implementar todos os padrões táticos do \acrshort{ddd} (como agregados explícitos ou eventos de domínio em toda a base), mas sim dar nome e limite às partes do sistema que já existiam no modelo de dados e nos fluxos de curadoria e de contribuição pública. A lista dos domínios adotados e a correspondência com modelos, controladores e serviços são detalhadas na Seção~\ref{sec:dev-ddd}.

% Seções previstas para expansão na monografia (fundamentação teórica, não resultados):
%\section{Testes de software}\label{sec:ref-testes}
%\section{Padrões de projeto e arquitetura de aplicações web}\label{sec:ref-padroes}

